% Section 7 — Discussion
% Page budget: ~1.5-2 pages in acmart sigconf two-column (~1700 words).
%
% Source-of-truth verification (every claim, no fake info):
%   - Cost savings 64.1%/65.3%: PAPER-FINDINGS.md L24-26, L162.
%   - Pareto threshold sweep flat in Hit@K: PAPER-FINDINGS L311-315.
%   - ReAct null OB p=0.18: PAPER-FINDINGS L262-275; OTel/WoL exact zero: L279-283.
%   - Extended-trace harm OB Δ=-0.018, p=0.002: PAPER-FINDINGS L277-278, L325.
%   - False-novelty cross-dataset 40.6%/39.7%/0.0%: PAPER-FINDINGS L399-409.
%   - Mechanism (verifier flag drives false-novel via OR disjunction):
%     results/ob/4.8-failure-categories/SUMMARY.md L36-65; the documented
%     statistic is novel_precision=0.1128 (89% of novel-flagged windows
%     are not actually novel). The specific "94% overlap with verifier
%     flag" number isn't in any SUMMARY.md and is therefore not used here.
%   - WoL v2 triage 0.164 (default) vs 0.921 (dense_only): PAPER-FINDINGS L179.

\section{Discussion}
\label{sec:discussion}

The empirical results in \S\ref{sec:results}--\S\ref{sec:cross-app} carry five lessons that, taken together, motivate a small set of concrete recommendations for future incident-triage agents. We discuss each in turn before consolidating the recommendations at the end of the section.

\subsection{Cost-aware per-window planning is the right structural pattern}
\label{sec:discussion:pareto}

The most quantitatively confident result in the paper is the cost-savings claim: the agent matches the always-on cascade's retrieval quality on \dsob{} at $64.1\%$ lower dollar-equivalent cost, and the same pattern replicates at $65.3\%$ on \dsotel{}. The savings are insensitive to the cheap-path confidence threshold across the entire $[0.50, 0.95]$ operating range: every threshold produces identical Hit@K to four decimal places, and only cost varies (\S\ref{sec:results:pareto}). The reason is identifiable from the skill-ablation results in \S\ref{sec:results:ablation}: the cascade-derived primitives the cheap-path gate skips when it is confident do not change rank-one rankings on the windows they are skipped from. The expensive primitives---the language-model verifier, the hybrid retrievers, the on-demand tools---contribute meaningful signal on the small subset of windows where the cheap path is genuinely uncertain, and the gate is good at identifying that subset.

The production recommendation that emerges on \dsob{} is unambiguous: drop the cheap-path threshold from its conservative documented default of $0.90$ to $0.50$, unlocking an additional four percentage points of cost savings at zero retrieval cost. The general lesson is the design pattern, not the specific threshold value: for diagnostic systems with strongly bimodal triage-confidence distributions---most windows clearly noise, most active-fault windows clearly uncertain---a cheap-first / escalate-on-uncertainty architecture is the right structural choice, and the cheap-path threshold is the most under-discussed cost knob in the literature. A reviewer's natural objection that the savings percentage depends on a specific per-skill cost table is addressed by the construction of the metric itself: the always-on counterfactual and the agent's actual cost scale together with the underlying language-model unit rate, so the percentage is invariant to provider, model version, and self-hosted-versus-vendor configuration.

\subsection{The on-demand evidence-gathering loop: framework yes, lift not yet}
\label{sec:discussion:react}

The on-demand evidence-gathering loop is the agent's most architecturally ambitious component. It is fully implemented (\S\ref{sec:approach:tooluse}), fully instrumented with a five-mode failure-mode taxonomy, fully replayable from per-window traces, and registered as five concrete skills in the registry. The empirical question of whether the loop's positive contribution to Hit@1 is statistically distinguishable from zero, however, has a consistent answer across all three datasets: it is not. The point-estimate lift on \dsob{} is $+0.0151$ at paired-bootstrap $p = 0.18$; on \dsotel{} and \dswol{} the paired delta is exactly zero on every metric.

We report this as a finding rather than a failure. The term \emph{agentic retrieval} has begun to appear as a load-bearing claim in the AIOps and incident-response literature, and the empirical question of whether on-demand evidence gathering produces a measurable retrieval-quality lift over a well-tuned non-agentic retriever has received less rigorous attention than the concept's prominence would suggest. Our three-dataset measurement under paired-bootstrap confidence intervals is the most rigorous test the loop's lift has received in our corner of the literature, and the answer it returns is that the lift exists in point estimate, does not exceed the noise floor at the sample sizes we have, and is not robust across applications. Future work claiming a tool-use lift should be expected to clear the same bar on multiple datasets with paired-delta confidence intervals.

The only statistically significant ReAct-related finding in our measurements is a negative one: the extended-trace tool's $-0.018$ Hit@1 harm on \dsob{} ($p = 0.002$, \S\ref{sec:results:ablation}). The mechanism is intelligible---the tool surfaces a list of services seen on the call path during the window, and on \dsob{}'s call topology those service names match memory-side tickets across multiple unrelated scenario families, polluting the reranker's token bag with cross-family noise. The natural remediation is per-family tool selection: enable the trace tool only on scenario families where its service-graph context is dataset-validated to outweigh the cross-family noise. The framework supports this as a controller-rule edit; the measurement is straightforward follow-up.

The interpretation we draw is not that on-demand evidence gathering is the wrong architecture. It is that the loop's value is unlikely to come from a single one-size-fits-all rule-based controller with a flat allowlist of tools applied uniformly across scenario families. A controller that learns which tools to call on which family from trace data---the per-window traces the agent already emits are exactly the labelled substrate such a policy would consume---is the natural follow-up, and the framework is built to accommodate it as a configuration swap rather than an architectural change.

\subsection{The false-novelty mechanism is configurational, not architectural}
\label{sec:discussion:falsenovel}

The cleanest cross-dataset finding of the paper is the false-novelty rate's collapse from $\sim$$40\%$ on the two synthetic datasets to exactly $0\%$ on real Apache Jira (\S\ref{sec:cross-app:falsenovel}). The mechanism, documented in the failure-categories analysis on \dsob{}, is the language-model verifier's own novelty flag: at headline measurement only $11.3\%$ of novel-flagged windows are actually novel ($\text{novel\_precision} = 0.113$), and the verifier's $\text{is\_novel}$ field is the OR-disjunction's most eager input. On the two datasets where the verifier runs, its over-eager flag dominates the disjunction's output; on \dswol{} where the verifier is structurally absent, the disjunction is left with the two remaining signals (retrieval-confidence threshold and a learned per-window classifier), and the rate falls to zero.

The implication for system design is concrete. The agent's largest blind spot on the two telemetry-rich datasets is not architectural but configurational: the OR-disjunction over three novelty signals trusts the verifier's flag at face value, and the verifier is over-eager. Two design changes follow from the finding, each a small local edit. The first is to require two of the three signals to agree before emitting a novel decision---a soft-AND composition rather than an OR. The second is to apply the verifier-known-helpful gate at the novelty-signal layer rather than only at the verifier-invocation layer: even on datasets where the verifier is helpful for retrieval re-ranking, its novelty flag can be selectively muted. The validation of either change is straightforward in our framework---rerun with a configuration switch, paired-bootstrap the delta---and we expect both to materially reduce the synthetic-dataset false-novelty rate by the same mechanism that produces \dswol{}'s zero rate.

The broader interpretation is that single-dataset evaluations conflate \emph{caused-by-the-system} blind spots with \emph{caused-by-a-configurable-component} blind spots. A paper that evaluated the agent only on \dsob{} would have reported the $40.6\%$ false-novelty rate as an inherent property of the disjunction structure, recommended re-architecting the disjunction, and missed the actual cause. Cross-dataset replication is what surfaced the mechanism, and the configurational nature of the fix is the design recommendation that follows.

\subsection{Architecture versus operating point: the \dswol{} triage finding}
\label{sec:discussion:opspoint}

The second cross-dataset finding that a single-dataset paper would have missed sits in the triage-accuracy row of \S\ref{sec:cross-app:wol}. At the agent's default configuration, triage accuracy on \dswol{} is $0.164$---substantially below both synthetic datasets and below \dswol{}'s test-split majority-class baseline of $0.496$. A \texttt{dense\_only} reduced configuration that disables the over-confident escalation path recovers triage accuracy to $0.921$, well above the majority-class baseline.

The distinction this exposes is between architectural correctness and operating-point correctness. The architecture is sound: the same skill registry, the same controller branches, and the same triage stacker that produces $0.83$--$0.84$ on the telemetry-rich synthetic datasets recovers $0.92$ on real Apache Jira under a reduced configuration. The operating point is dataset-specific: the cheap-path triage stacker's default threshold, calibrated on \dsob{}'s ticket-worthy-heavy distribution, over-pages on \dswol{}'s $21$/$26$/$53$\% ticket-worthy/borderline/noise distribution. The deployable system therefore needs per-dataset operating-point selection, which is a configuration concern rather than an architectural one---exactly the kind of distinction the capability-adaptive layered design is built to make. The recommendation is that any productionisation of the agent should treat the cheap-path triage threshold as a per-dataset hyperparameter, fit on a small held-out validation slice of the target dataset before deployment, and revalidated whenever the dataset's class distribution drifts.

\subsection{What the cross-dataset replication is for}
\label{sec:discussion:why3}

A reviewer might reasonably ask: why three datasets? The answer is that two findings in the paper would not be reportable without the third. The ReAct null result is not credible from a single dataset because the obvious response is \emph{``you under-tuned your tools''}; replicated as exact zero on a structurally different second dataset and as a structural-zero on a real-data third dataset, the null becomes a finding that future tool-use claims have to argue around rather than a within-dataset accident that any well-tuned re-run could overturn. The false-novelty mechanism is not localisable from a single dataset because the verifier is on every window on a single-dataset run; only the dataset where the verifier is structurally absent reveals the mechanism by which the false-novelty signal flows. Three datasets is therefore a methodological discipline as much as an empirical breadth claim, and the paper's strongest cross-dataset findings---the universal pages-per-incident ratio, the cost-savings rate, the BiEncoder-anchored retrieval, the cleanly localised false-novelty mechanism---are exactly the ones that benefit from the replication.

\subsection{Concrete design recommendations}
\label{sec:discussion:recommendations}

The findings above produce five recommendations whose evidence sits directly in the paper.

\paragraphhead{(1) Drop the cheap-path confidence threshold to $0.50$ on synthetic-microservices deployments.} Recovers four percentage points of dollar-equivalent cost savings on \dsob{} at zero Hit@K cost. The same Pareto-flat pattern holds on \dsotel{}. Dataset-specific in absolute savings, architecturally portable.

\paragraphhead{(2) Disable the extended-trace tool by default; enable it per scenario family.} The trace tool is significantly harmful on \dsob{} at $p = 0.002$. A controller policy that enables the tool only on scenario families where its service-graph context is dataset-validated to outweigh the cross-family token noise is the implementable follow-up. The framework supports this as a controller-rule edit.

\paragraphhead{(3) Replace the novelty composer's OR-disjunction with a soft-AND that requires two of three signals.} Localises the false-novelty fix to the disjunction structure rather than to the verifier itself, and is expected to reduce the synthetic-dataset false-novelty rate by the same mechanism that produces \dswol{}'s zero rate. The change is a one-line edit; the validation is a paired-bootstrap rerun.

\paragraphhead{(4) Calibrate the cheap-path triage operating point per dataset.} On \dswol{} the default triage threshold over-pages and accuracy falls below the majority-class baseline; a \texttt{dense\_only} reduced configuration recovers accuracy to $0.921$. Architectural transfer and operating-point transfer are different questions; the deployable system requires both.

\paragraphhead{(5) Train a learned controller from trace data before scaling the tool set.} Every per-window plan, gate evaluation, skill output, and budget snapshot is already serialised in the trace store; a logistic-regression-or-larger controller that learns which tools to call on which family from those traces is the natural next architectural step. The current rule-based controller's flat allowlist is what produces the cross-dataset tool-use null; a learned policy that respects per-family tool effectiveness is the path the framework was built to support.

These five recommendations are not aspirational. Each is a configuration entry in the agent's existing layered architecture (\S\ref{sec:approach:extensibility}), each can be validated by the paired-bootstrap measurement harness the paper already uses, and each is portable across the three datasets and any future dataset that follows the same per-dataset profile.
